\usepackage{cite}
\usepackage{graphicx}
\graphicspath{{./} {fig/}}
% \usepackage[tight,footnotesize]{subfigure}
\usepackage[labelsep=period]{caption}
\captionsetup[figure]{font=small}
\captionsetup[table]{textfont={sc,footnotesize}, labelfont=footnotesize, labelsep=newline, justification=centering} %ieee table format
\usepackage{subcaption}
\renewcommand\thesubfigure{(\alph{subfigure})} %引用子图格式：数字(字母)，如Fig.1(a)
\DeclareCaptionLabelFormat{opening}{#2}
\captionsetup[subfigure]{labelformat=opening} %避免子图出现双括号

%\usepackage{amsmath}
\usepackage{amsmath, amssymb} % 强烈建议
\usepackage{algorithmic}
% \usepackage{url}
\usepackage[hyphens,spaces,obeyspaces]{url}
\usepackage{footmisc}
\def\myurl#1{\setbox0\vbox{\hsize.5\maxdimen
\url{#1}\par
\global\setbox1\lastbox}\unhbox1 }
\usepackage[colorlinks]{hyperref}
\usepackage{array}

\usepackage{threeparttable}
\usepackage{tabularx}
\usepackage{booktabs}
\usepackage[table,xcdraw]{xcolor}
\usepackage{adjustbox}
\usepackage[utf8]{inputenc}
\usepackage{stfloats}
\usepackage{array} % 记得在导言区加这个包，否则 m{} 无法使用
\usepackage{acronym}
\usepackage{multirow}
\usepackage{algorithm}
\usepackage{algorithmic}
\newcommand{\chicca}[1]{{\color{red}{\bf elisabetta:} #1}}

%% Block comment. \begin{comment}  \end{comment}
\usepackage{verbatim}

%% Provides generic commands \degree, \celsius, \perthousand, \micro and \ohm which work both in text and maths mode.
\usepackage{textcomp,gensymb}

%% line break
\usepackage{makecell}

%% 调整表格线与上下内容的间隔
\usepackage{booktabs}

\usepackage{hyperref}
\hypersetup{
   hypertex=true,
   colorlinks=true,
   linkcolor=blue,
   anchorcolor=blue,
   citecolor=blue
}

%% emphasizing text 
\usepackage{soul}
\setstcolor{red}

\newsavebox\CBox
\def\textBF#1{\sbox\CBox{#1}\resizebox{\wd\CBox}{\ht\CBox}{\textbf{#1}}}
